% Generated by publication/build_sections.py; do not edit by hand.
\paragraph{Informal verdict.}
\textbf{A --- Complete solution}, with high confidence.

\paragraph{Lean coverage (partial).}
Lean proves the exact Hadamard integral representation, the dominant-saddle extremal lemma, the sharp exponential envelope for every fixed m, and the complete Bessel identity for m=2. The general oscillatory Poincare expansion, algebraic prefactor, correction coefficient, and uniform remainder remain under Aristotle continuation.

\subsection{Audited informal answer}
Q776 --- Restatement
For each fixed integer \(m\ge 2\), define the entire function

\[
f_m(x)=\sum_{n=0}^{\infty}\frac{x^n}{(n!)^m}.
\]

The problem asks for a simple asymptotic equivalent of \(f_m(x)\) as \(x\to-\infty\). The editorial remark also records the positive-axis asymptotic

\[
f_m(x)\sim
(2\pi)^{(1-m)/2}m^{-1/2}
x^{(1-m)/(2m)}e^{m x^{1/m}}
\qquad (x\to+\infty).
\]

This is the complete statement on page 14 of the canonical PDF. Luc Pommeret
There is no apparent typo. There is, however, a genuine ambiguity in the word ``equivalent'': on the negative axis the dominant term is oscillatory and has infinitely many zeros, so it cannot be a global equivalent in the strict ratio sense. The correct global result is an additive Poincaré asymptotic expansion.

Final theorem
Fix \(m\ge2\). Put

\[
t=-x,\qquad R=t^{1/m},\qquad
\alpha=\frac{\pi}{m},
\]

and define

\[
\Psi_m(R)
=
mR\sin\alpha-\frac{m-1}{2}\alpha
=
mR\sin\frac{\pi}{m}
-\frac{(m-1)\pi}{2m}.
\]

Then, as \(x\to-\infty\), equivalently \(R\to+\infty\),

\[
\boxed{
\begin{aligned}
f_m(x)
={}&
\frac{2(2\pi)^{(1-m)/2}}{\sqrt m}\,
t^{(1-m)/(2m)}
e^{m t^{1/m}\cos(\pi/m)}
\\
&\times
\left[
\cos\Psi_m(t^{1/m})
+
\frac{m^2-1}{24m\,t^{1/m}}
\cos\left(
\Psi_m(t^{1/m})-\frac{\pi}{m}
\right)
+
O_m\!\left(t^{-2/m}\right)
\right].
\end{aligned}}
\]

In particular, the leading global estimate is

\[
\boxed{
f_m(x)
=
\frac{2(2\pi)^{(1-m)/2}}{\sqrt m}\,
(-x)^{(1-m)/(2m)}
e^{m(-x)^{1/m}\cos(\pi/m)}
\left[
\cos\!\left(
m(-x)^{1/m}\sin\frac{\pi}{m}
-\frac{(m-1)\pi}{2m}
\right)
+
O_m\!\left((-x)^{-1/m}\right)
\right].
}
\]

The constants implicit in the \(O\)-terms depend on the fixed integer \(m\), but not on \(x\).
More generally, there are real constants \(c_{m,k}\), with

\[
c_{m,0}=1,
\qquad
c_{m,1}=\frac{m^2-1}{24m},
\]

such that, for every fixed \(N\ge1\),

\[
\boxed{
\begin{aligned}
f_m(-R^m)
={}&
\frac{2(2\pi)^{(1-m)/2}}{\sqrt m}\,
R^{-(m-1)/2}e^{mR\cos(\pi/m)}
\\
&\times
\left[
\sum_{k=0}^{N-1}
\frac{c_{m,k}}{R^k}
\cos\left(
\Psi_m(R)-\frac{k\pi}{m}
\right)
+
O_{m,N}(R^{-N})
\right].
\end{aligned}}
\]

This is a full asymptotic expansion in descending integral powers of \(R=(-x)^{1/m}\).

Proof
\begin{enumerate}
\item An exact torus integral representation
\end{enumerate}
Let \(d=m-1\), and for
\(\theta=(\theta_1,\ldots,\theta_d)\in[-\pi,\pi]^d\), set

\[
S(\theta)
=
\sum_{j=1}^{d}e^{i\theta_j}
+
e^{-i(\theta_1+\cdots+\theta_d)}.
\]

For every \(r\in\mathbb C\),

\[
\boxed{
f_m(r^m)
=
\frac{1}{(2\pi)^d}
\int_{[-\pi,\pi]^d}
e^{rS(\theta)}\,d\theta.
}
\]

Indeed, expanding all exponentials gives

\[
e^{rS(\theta)}
=
\prod_{j=1}^{d}
\left(
\sum_{n_j\ge0}
\frac{r^{n_j}e^{in_j\theta_j}}{n_j!}
\right)
\left(
\sum_{n_m\ge0}
\frac{r^{n_m}e^{-in_m(\theta_1+\cdots+\theta_d)}}{n_m!}
\right).
\]

The expansion is absolutely and uniformly convergent on the compact torus. Integration with respect to each \(\theta_j\) forces

\[
n_j=n_m\qquad(1\le j\le d).
\]

Consequently,

\[
\frac{1}{(2\pi)^d}\int e^{rS(\theta)}\,d\theta
=
\sum_{n=0}^{\infty}\frac{r^{mn}}{(n!)^m}
=
f_m(r^m).
\]

For the negative axis choose

\[
r=Re^{i\alpha},
\qquad
\alpha=\frac{\pi}{m},
\]

so that \(r^m=-R^m\).

\begin{enumerate}
\item Identification of the dominant saddle points for \(m\ge3\)
\end{enumerate}
We first prove an elementary extremal lemma.
Lemma
Let \(z_1,\ldots,z_m\) be complex numbers of modulus \(1\) satisfying

\[
z_1z_2\cdots z_m=-1.
\]

For \(m\ge3\),

\[
\operatorname{Re}(z_1+\cdots+z_m)
\le m\cos\frac{\pi}{m}.
\]

Equality holds exactly in the following two cases:

\[
z_1=\cdots=z_m=e^{i\pi/m},
\]

or

\[
z_1=\cdots=z_m=e^{-i\pi/m}.
\]

Proof
The constraint set is compact, so a maximum exists. Write \(z_j=e^{iu_j}\). Tangent variations preserving the product have the form

\[
u_j\longmapsto u_j+s t_j,
\qquad
\sum_{j=1}^{m}t_j=0.
\]

At a critical point,

\[
\sum_{j=1}^{m}(-\sin u_j)t_j=0
\]

for every \((t_j)\) of sum zero. Hence all the numbers \(\sin u_j\) are equal.
Thus, for some \(\beta\), each \(u_j\) is congruent modulo \(2\pi\) either to \(\beta\) or to \(\pi-\beta\). Suppose \(p\) of the angles are of the first type and \(m-p\) are of the second type. Then

\[
\operatorname{Re}\sum_{j=1}^{m}z_j
=
(2p-m)\cos\beta.
\]

If \(0<p<m\), then

\[
\operatorname{Re}\sum z_j\le |2p-m|\le m-2.
\]

On the other hand,

\[
m-m\cos\frac{\pi}{m}
=
2m\sin^2\frac{\pi}{2m}
<
\frac{\pi^2}{2m}
<2
\qquad(m\ge3),
\]

so

\[
m-2<m\cos\frac{\pi}{m}.
\]

Therefore no mixed configuration can maximize the real part.
At a maximizer all the \(z_j\) are equal, say \(z_j=z\). The product condition gives \(z^m=-1\). Among the \(m\)-th roots of \(-1\), the largest real part is \(\cos(\pi/m)\), attained exactly by \(e^{\pm i\pi/m}\). \(\square\)
Apply the lemma with

\[
z_j=e^{i(\alpha+\theta_j)}
\quad(1\le j\le d),
\qquad
z_m=e^{i(\alpha-\theta_1-\cdots-\theta_d)}.
\]

Their product is

\[
z_1\cdots z_m=e^{im\alpha}=-1,
\]

and

\[
\operatorname{Re}(rS(\theta))
=
R\,\operatorname{Re}(z_1+\cdots+z_m).
\]

Thus the maximum of the real part of the exponent is

\[
mR\cos\alpha,
\]

and there are exactly two maximizers on the torus:

\[
\theta^{(+)}=(0,\ldots,0)
\]

and

\[
\theta^{(-)}=(-2\alpha,\ldots,-2\alpha).
\]

Outside fixed disjoint neighborhoods of these points, compactness gives some
\(\eta_m>0\) such that

\[
\operatorname{Re}(rS(\theta))
\le
mR\cos\alpha-\eta_mR.
\]

The corresponding part of the integral is therefore

\[
O_m\!\left(e^{mR\cos\alpha-\eta_mR}\right),
\]

which is exponentially smaller than every term retained below.

\begin{enumerate}
\item Local expansion at a saddle
\end{enumerate}
We analyze the saddle at the origin. Introduce

\[
v_j=u_j\quad(1\le j\le d),
\qquad
v_m=-\sum_{j=1}^{d}u_j.
\]

Then \(\sum_{j=1}^{m}v_j=0\), and locally

\[
S(u)=\sum_{j=1}^{m}e^{iv_j}.
\]

Taylor expansion gives

\[
S(u)
=
m-\frac12Q(u)+H_3(u)+H_4(u)+H_5(u)+O_m(|u|^6),
\]

where

\[
Q(u)=\sum_{j=1}^{m}v_j^2,
\]


\[
H_3(u)=-\frac{i}{6}\sum_{j=1}^{m}v_j^3,
\qquad
H_4(u)=\frac1{24}\sum_{j=1}^{m}v_j^4,
\]

and \(H_5\) is homogeneous and odd.
The quadratic form is

\[
Q(u)
=
\sum_{j=1}^{d}u_j^2
+
\left(\sum_{j=1}^{d}u_j\right)^2
=
u^{\mathsf T}Au,
\]

with

\[
A=I_d+\mathbf 1\mathbf 1^{\mathsf T}.
\]

Its eigenvalues are \(1\), with multiplicity \(d-1\), and \(m\), with multiplicity \(1\). Hence

\[
\det A=m.
\]

Let \(\rho\to\infty\) in a fixed closed subsector of \(\operatorname{Re}\rho>0\). A standard scaling \(u=|\rho|^{-1/2}y\), made rigorous by splitting into
\(|y|\le |\rho|^\varepsilon\) and its complement, yields

\[
\begin{aligned}
\int_{\text{near }0}e^{\rho(S(u)-m)}\,du
={}&
\frac{(2\pi)^{d/2}}{\sqrt m}\,
\rho^{-d/2}
\\
&\times
\left(
1+\frac{c_{m,1}}{\rho}
+O_m(|\rho|^{-2})
\right).
\end{aligned}
\]

For completeness, we calculate \(c_{m,1}\).
After the Gaussian scaling, the potential contribution of order
\(\rho^{-1/2}\) is odd and integrates to zero. The contribution of order
\(\rho^{-1}\) is

\[
H_4+\frac12H_3^2.
\]

Let \(V=(V_1,\ldots,V_m)\) be the centered Gaussian vector on the hyperplane
\(\sum V_j=0\) corresponding to the density \(e^{-Q/2}\). Its covariance matrix is

\[
\mathbb E(V_jV_k)=\delta_{jk}-\frac1m.
\]

In particular,

\[
v:=\mathbb E(V_j^2)=\frac{m-1}{m},
\qquad
c:=\mathbb E(V_jV_k)=-\frac1m
\quad(j\ne k).
\]

By Gaussian pairings,

\[
\mathbb E(V_j^4)=3v^2,
\]

so

\[
\mathbb E\sum_{j=1}^{m}V_j^4
=
\frac{3(m-1)^2}{m}.
\]

Also,

\[
\mathbb E(V_j^6)=15v^3,
\]

and, for \(j\ne k\),

\[
\mathbb E(V_j^3V_k^3)
=
9v^2c+6c^3.
\]

Therefore

\[
\begin{aligned}
\mathbb E\left(\sum_{j=1}^{m}V_j^3\right)^2
&=
15m v^3
+
m(m-1)(9v^2c+6c^3)
\\
&=
\frac{6(m-1)(m-2)}{m}.
\end{aligned}
\]

It follows that

\[
\begin{aligned}
c_{m,1}
&=
\frac1{24}
\mathbb E\sum_jV_j^4
-
\frac1{72}
\mathbb E\left(\sum_jV_j^3\right)^2
\\
&=
\frac{(m-1)^2}{8m}
-
\frac{(m-1)(m-2)}{12m}
\\
&=
\boxed{\frac{m^2-1}{24m}}.
\end{aligned}
\]

The same scaling argument, carried to arbitrary order, gives

\[
\int_{\text{near }0}e^{\rho(S(u)-m)}\,du
\sim
\frac{(2\pi)^{d/2}}{\sqrt m}\,
\rho^{-d/2}
\sum_{k=0}^{\infty}\frac{c_{m,k}}{\rho^k}.
\]

Only integral powers occur: at every half-integral order the polynomial in the scaled variables is odd, so its Gaussian integral vanishes.

\begin{enumerate}
\item The two conjugate saddle contributions for \(m\ge3\)
\end{enumerate}
At the first saddle take

\[
\rho=r=Re^{i\alpha}.
\]

Its contribution is

\[
\frac{(2\pi)^{-d/2}}{\sqrt m}\,
e^{mr}r^{-d/2}
\left(
1+\frac{c_{m,1}}{r}+O_m(R^{-2})
\right).
\]

For the second saddle, write

\[
\theta=\theta^{(-)}+u.
\]

An exact identity holds:

\[
S(\theta^{(-)}+u)=e^{-2i\alpha}S(u).
\]

Indeed, every one of the \(m\) exponential terms acquires the same factor
\(e^{-2i\alpha}\). Consequently,

\[
rS(\theta^{(-)}+u)
=
\overline r\,S(u),
\]

because

\[
re^{-2i\alpha}=Re^{-i\alpha}=\overline r.
\]

Thus the second saddle contributes the complex conjugate

\[
\frac{(2\pi)^{-d/2}}{\sqrt m}\,
e^{m\overline r}\overline r^{-d/2}
\left(
1+\frac{c_{m,1}}{\overline r}+O_m(R^{-2})
\right).
\]

Adding the two contributions and absorbing the exponentially smaller remainder gives

\[
\begin{aligned}
f_m(-R^m)
={}&
\frac{(2\pi)^{-d/2}}{\sqrt m}
\bigg[
e^{mr}r^{-d/2}
\left(1+\frac{c_{m,1}}r+O_m(R^{-2})\right)
\\
&\hspace{35mm}
+
e^{m\bar r}\bar r^{-d/2}
\left(1+\frac{c_{m,1}}{\bar r}+O_m(R^{-2})\right)
\bigg].
\end{aligned}
\]

Now

\[
e^{mr}r^{-d/2}
=
R^{-d/2}e^{mR\cos\alpha}
e^{\,i(mR\sin\alpha-d\alpha/2)}.
\]

Therefore

\[
\begin{aligned}
f_m(-R^m)
={}&
\frac{2(2\pi)^{-d/2}}{\sqrt m}
R^{-d/2}e^{mR\cos\alpha}
\\
&\times
\left[
\cos\left(mR\sin\alpha-\frac d2\alpha\right)
+
\frac{c_{m,1}}R
\cos\left(mR\sin\alpha-\frac d2\alpha-\alpha\right)
+
O_m(R^{-2})
\right].
\end{aligned}
\]

Since \(d=m-1\), this is the claimed formula for \(m\ge3\).

\begin{enumerate}
\item The exceptional boundary case \(m=2\)
\end{enumerate}
Here \(\alpha=\pi/2\), so \(\operatorname{Re}r=0\); the real-part maximum used above is no longer isolated. A stationary-phase argument is required.
The integral representation becomes

\[
f_2(-R^2)
=
\frac1{2\pi}\int_{-\pi}^{\pi}e^{2iR\cos\theta}\,d\theta.
\]

Equivalently,

\[
f_2(-R^2)=J_0(2R).
\]

The stationary points are \(\theta=0\) and \(\theta=\pi\).
Near \(\theta=0\), make the exact substitution

\[
u=2\sin\frac{\theta}{2}.
\]

Then

\[
2\cos\theta=2-u^2,
\qquad
d\theta=\frac{du}{\sqrt{1-u^2/4}}.
\]

With a smooth cutoff equal to \(1\) near the saddle, the local amplitude \(a(u)\) satisfies

\[
a(0)=1,
\qquad
a''(0)=\frac14.
\]

For a smooth compactly supported amplitude,

\[
\int_{\mathbb R}e^{-iRu^2}a(u)\,du
=
e^{-i\pi/4}\sqrt{\frac{\pi}{R}}
\left[
a(0)-\frac{i}{4R}a''(0)+O(R^{-2})
\right].
\]

One direct proof uses the Fourier transform and the exact Fresnel identity

\[
\int_{\mathbb R}e^{-iRu^2+i\xi u}\,du
=
e^{-i\pi/4}\sqrt{\frac{\pi}{R}}
e^{i\xi^2/(4R)}.
\]

Thus the contribution of \(\theta=0\) is

\[
\frac{e^{i(2R-\pi/4)}}{2\sqrt{\pi R}}
\left(
1-\frac{i}{16R}+O(R^{-2})
\right).
\]

The contribution of \(\theta=\pi\) is its complex conjugate. Away from the two stationary points, repeated integration by parts gives an error smaller than any prescribed power of \(R^{-1}\). Hence

\[
\boxed{
f_2(-R^2)
=
\frac1{\sqrt{\pi R}}
\left[
\cos\left(2R-\frac{\pi}{4}\right)
+
\frac1{16R}
\sin\left(2R-\frac{\pi}{4}\right)
+
O(R^{-2})
\right].
}
\]

This is exactly the general formula, because

\[
\frac{m^2-1}{24m}\bigg|_{m=2}=\frac1{16},
\]

and

\[
\cos\left(\Psi_2(R)-\frac{\pi}{2}\right)=\sin\Psi_2(R).
\]

Thus all \(m\ge2\) are covered.

Why the leading term is not a global strict equivalent
Let

\[
\mathcal A_m(R)
=
\frac{2(2\pi)^{(1-m)/2}}{\sqrt m}
R^{-(m-1)/2}e^{mR\cos(\pi/m)}.
\]

The leading approximation would be

\[
\mathcal A_m(R)\cos\Psi_m(R).
\]

For integers \(k\) large enough, define

\[
R_k^*
=
\frac{
\frac{\pi}{2}+k\pi+\frac{(m-1)\pi}{2m}
}{
m\sin(\pi/m)
}.
\]

Then

\[
\Psi_m(R_k^*)=\frac{\pi}{2}+k\pi,
\]

so the proposed leading expression is exactly zero. But the two-term expansion gives

\[
\begin{aligned}
f_m(-(R_k^*)^m)
&=
\mathcal A_m(R_k^*)
\left[
\frac{m^2-1}{24mR_k^*}
\cos\left(
\frac{\pi}{2}+k\pi-\frac{\pi}{m}
\right)
+
O_m((R_k^*)^{-2})
\right]
\\
&=
(-1)^k\mathcal A_m(R_k^*)
\left[
\frac{m^2-1}{24mR_k^*}
\sin\frac{\pi}{m}
+
O_m((R_k^*)^{-2})
\right],
\end{aligned}
\]

which is nonzero for all sufficiently large \(k\).
Therefore

\[
f_m(-R^m)\sim
\mathcal A_m(R)\cos\Psi_m(R)
\]

is not valid globally in the strict quotient sense.
It is valid along every approach for which

\[
R\,|\cos\Psi_m(R)|\longrightarrow\infty.
\]

In particular, it holds uniformly on sets satisfying

\[
|\cos\Psi_m(R)|\ge\delta>0.
\]

The additive expansion stated in the theorem is valid without excluding any values of \(R\).

Checks and boundary behavior
\(m=2\)

\[
f_2(-t)=J_0(2\sqrt t),
\]

and the theorem gives

\[
f_2(-t)
=
\frac{t^{-1/4}}{\sqrt\pi}
\left[
\cos\left(2\sqrt t-\frac{\pi}{4}\right)
+
\frac{1}{16\sqrt t}
\sin\left(2\sqrt t-\frac{\pi}{4}\right)
+
O(t^{-1})
\right],
\]

the classical Bessel expansion.
\(m=3\)
Since \(\cos(\pi/3)=1/2\) and \(\sin(\pi/3)=\sqrt3/2\),

\[
\begin{aligned}
f_3(-t)
={}&
\frac{e^{\frac32t^{1/3}}}{\pi\sqrt3\,t^{1/3}}
\bigg[
\cos\left(
\frac{3\sqrt3}{2}t^{1/3}-\frac{\pi}{3}
\right)
\\
&\qquad+
\frac1{9t^{1/3}}
\cos\left(
\frac{3\sqrt3}{2}t^{1/3}-\frac{2\pi}{3}
\right)
+
O(t^{-2/3})
\bigg].
\end{aligned}
\]

Thus \(f_3(-t)\) oscillates with an exponentially increasing envelope.
Sign changes
At points where \(\Psi_m(R)=k\pi\),

\[
f_m(-R^m)
=
(-1)^k\mathcal A_m(R)(1+O_m(R^{-1})).
\]

Hence \(f_m\) changes sign infinitely often on the negative real axis, for every \(m\ge2\).
Necessity of the hypotheses
The integer condition on \(m\) is essential to the elementary \(m\)-torus representation used in the proof. The case \(m=1\), excluded in the statement, is qualitatively different:

\[
f_1(x)=e^x,
\qquad
f_1(x)=e^x\quad(x\to-\infty).
\]

The estimates above are for fixed \(m\). They are not uniform as \(m\to\infty\): the two dominant saddles are separated by \(2\pi/m\) and coalesce in that additional limit.
